\documentclass[10pt,a4paper]{amsart}
\usepackage[margin=19mm]{geometry}
\usepackage{amsmath,amssymb,mathtools}
\usepackage[scr=rsfs,cal=euler]{mathalfa}
\usepackage{xfrac}
\usepackage[unicode,hidelinks,bookmarks=false]{hyperref}
\newcommand{\ie}{i.e.\ }
\newcommand{\cf}{cf.\ }
\newcommand{\resp}{respectively\ }
\newcommand{\Subsection}{\subsection}
\providecommand{\nobreakdash}{\nobreak\hbox{-}}
\setlength{\emergencystretch}{2em}
\multlinegap0pt
\usepackage{amssymb}
\usepackage[shortlabels]{enumitem}
\newtheorem{theorem}{Theorem}
\title{A Generalized Cross Ratio}
\author{Michael R. Pilla}
\date{Source: arXiv:2106.10570v1 (CC BY 4.0)}
\begin{document}
\maketitle

\noindent\emph{Source and licence.} A Generalized Cross Ratio. Version arXiv:2106.10570v1 (CC BY 4.0), distributed by arXiv under the Creative Commons Attribution 4.0 International licence, http://creativecommons.org/licenses/by/4.0/, as recorded in the arXiv licence metadata for this exact version and shown on the abstract page https://arxiv.org/abs/2106.10570v1. Full licence notice: \texttt{licenses/arxiv-2106.10570-CC-BY-4.0.txt}.\par\smallskip

\noindent\textit{Editorial scope.} Counted unit: the article's theorem theorem (source0.tex lines 407--409). The article's complete original proof of it is delivered with it (source0.tex lines 411--491, one \texttt{proof} environment of 81 lines). No mathematics is written, repaired or re-derived: the statement and the proof are the article's own lines, in the article's own notation, and no retained line is substituted or re-flowed.  No further source-local context is needed: every cross-reference of every retained line resolves inside the two delivered spans.  Nothing was omitted from any retained span, so the package carries no omission record.  No retained line contains a citation, so the wrapper carries no bibliography and no bibliography entry.  No retained line contains a figure environment, an image-inclusion command or drawing code, so no image is delivered and the package declares no asset dependency.  Editorial formatting only, in addition: an AMS \texttt{amsart} preamble that restates the article's own declarations and macros used by the retained lines, namely the environments \texttt{theorem} on separate counters, together with the standard packages those lines need.  The retained environments print in this document as Theorem 1 in order of appearance, whereas the article prints the same items with the numbering of its own full document.  The complete native source of the article as distributed in the arXiv e-print for this exact version is bundled unchanged as \texttt{sources/112834/source0.tex}.
\begin{theorem}
Identifying coordinates $z^i_j=\frac{u^i_j}{u^{N+1}_j}$ for $j=1,...,N$ and $u^{N+3}=(0,0,...,0,1,0)^T$, the cross ratio in ${\bf CP}^N$ reduces to the cross ratio in ${\bf CP}^{N-1}$.
\end{theorem}

\begin{proof}
We use proof by induction. We have seen that our result is true for $N=2$. Now suppose $N=k+1$ where $k>1$, making the appropriate identifications, we have
\begin{align*}
&(u^1, u^2, ..., u^{(k+1)+3})=\frac{[u^1, u^3, u^5,..., u^{k+4}][u^2, u^4, u^5,..., u^{k+4}]}{[u^1, u^4, u^5,..., u^{k+4}][u^2, u^3, u^5,..., u^{k+4}]}\\
&=\frac{ \det \begin{pmatrix}
    u^1_1 & u^3_1 & u_1^5 & \dots  & u^{k+4}_1 \\
   u^1_2 & u^3_2 & u_2^5 & \dots  & u^{k+4}_2 \\
    \vdots & \vdots & \vdots & \ddots & \vdots \\
    u^1_{k+2} & u^3_{k+2} & u_{k+2}^5 & \dots  & u^{k+4}_{k+2}
\end{pmatrix}\det \begin{pmatrix}
    u_1^2 & u_1^4 & u_1^5 & \dots  & u_1^{k+4} \\
    u_2^2 & u_2^4 & u_2^5 & \dots  & u_2^{k+4} \\
    \vdots & \vdots & \vdots & \ddots & \vdots \\
    u_{k+2}^2 & u_{k+2}^4 & u_{k+2}^5 & \dots  & u_{k+2}^{k+4}
\end{pmatrix}}
{\det \begin{pmatrix}
    u_1^1 & u_1^4 & u_1^5 & \dots  & u_1^{k+4} \\
    u_2^1 & u_2^4 & u_2^5 & \dots  & u_2^{k+4} \\
    \vdots & \vdots & \vdots & \ddots & \vdots \\
    u_{k+2}^1 & u_{k+2}^4 & u_{k+2}^5 & \dots  & u_{k+2}^{k+4}
\end{pmatrix}\det \begin{pmatrix}
    u_1^2 & u_1^3 & u_1^5 & \dots  & u_1^{k+4} \\
    u_2^2 & u_2^3 & u_2^5 & \dots  & u_2^{k+4} \\
    \vdots & \vdots & \vdots & \ddots & \vdots \\
    u_{k+2}^2 & u_{k+2}^3 & u_{k+2}^5 & \dots  & u_{k+2}^{k+4}
\end{pmatrix}}\\
&=\frac{ \det \begin{pmatrix}
    z_1^1 & z_1^3 & z_1^5 & \dots  & 0 \\
    z_2^1 & z_2^3 & z_2^5 & \dots  & 0 \\
    \vdots & \vdots & \vdots & \ddots & \vdots \\
    z_{k+1}^1 & z_{k+1}^3 & z_{k+1}^5 & \dots  & 1 \\
    1 & 1 & 1& \dots  & 0
\end{pmatrix}\det \begin{pmatrix}
    z_1^2 & z_1^4 & z_1^5 & \dots  & 0 \\
    z_2^2 & z_2^4 & z_2^5 & \dots  & 0 \\
    \vdots & \vdots & \vdots & \ddots & \vdots \\
    z_{k+1}^2 & z_{k+1}^4 & z_{k+1}^5 & \dots  & 1 \\
    1 & 1 & 1& \dots  & 0
\end{pmatrix}}
{\det \begin{pmatrix}
    z_1^1 & z_1^4 & z_1^5 & \dots  & 0 \\
    z_2^1 & z_2^4 & z_2^5 & \dots  & 0 \\
    \vdots & \vdots & \vdots & \ddots & \vdots \\
    z_{k+1}^1 & z_{k+1}^4 & z_{k+1}^5 & \dots  & 1 \\
    1 & 1 & 1& \dots  & 0
\end{pmatrix}\det \begin{pmatrix}
    z_1^2 & z_1^3 & z_1^5 & \dots  & 0 \\
    z_2^2 & z_2^3 & z_2^5 & \dots  & 0 \\
    \vdots & \vdots & \vdots & \ddots & \vdots \\
    z_{k+1}^2 & z_{k+1}^3 & z_{k+1}^5 & \dots  & 1 \\
    1 & 1 & 1& \dots  & 0
\end{pmatrix}}\\ 
&=\frac{ \det \begin{pmatrix}
    z_1^1 & z_1^3 & z_1^5 & \dots  &  z_1^{k+3} \\
    z_2^1 & z_2^3 & z_2^5 & \dots  & z_2^{k+3}  \\
    \vdots & \vdots & \vdots & \ddots & \vdots \\
    z_{k}^1 & z_{k}^3 & z_{k}^5 & \dots  & z_{k}^{k+3}  \\
    1 & 1 & 1 & \dots & 1
\end{pmatrix}\det \begin{pmatrix}
    z_1^2 & z_1^4 & z_1^5 & \dots  & z_1^{k+3} \\
    z_2^2 & z_2^4 & z_2^5 & \dots  & z_2^{k+3} \\
    \vdots & \vdots & \vdots & \ddots & \vdots \\
    z_{k}^2 & z_{k}^4 & z_{k}^5 & \dots  &  z_{k}^{k+3} \\
    1 & 1 & 1 & \dots & 1
\end{pmatrix}}
{\det \begin{pmatrix}
    z_1^1 & z_1^4 & z_1^5 & \dots  & z_1^{k+3} \\
    z_2^1 & z_2^4 & z_2^5 & \dots  & z_2^{k+3} \\
    \vdots & \vdots & \vdots & \ddots & \vdots \\
    z_{k}^1 & z_{k}^4 & z_{k}^5 & \dots  &  z_{k}^{k+3} \\
    1 & 1 & 1 & \dots & 1
\end{pmatrix}\det \begin{pmatrix}
    z_1^2 & z_1^3 & z_1^5 & \dots  & z_1^{k+3} \\
    z_2^2 & z_2^3 & z_2^5 & \dots  & z_2^{k+3} \\
    \vdots & \vdots & \vdots & \ddots & \vdots \\
    z_{k}^2 & z_{k}^3 & z_{k}^5 & \dots  &  z_{k}^{k+3} \\
    1 & 1 & 1 & \dots & 1
\end{pmatrix}}\\
&=(z^1, z^2,..., z^{k+3}).
\end{align*}
\end{proof}

\end{document}
